% =====================================================================
% 2. YÖNTEM (Methodology)
% Evidence base: project_structure.md §2 (verified facts), configs/method/*.yaml,
%                configs/training/{frozen_extraction,finetune_wide}.yaml,
%                src/models/* docstrings, docs/decisions.md (PLD-04..08).
% Only IMPLEMENTED methods are described. No results here.
% =====================================================================
\section{Yöntem}
\label{sec:methodology}

Önerilen mimari üç bileşenden oluşur: önceden eğitilmiş CNN omurga ağları ile
özellik çıkarımı, bu özelliklerin ortak bir boyuta izdüşürülüp füzyonlanması ve
birleşik temsilin bir MLP sınıflandırıcı ile sınıflandırılması. Bu bölümde her
bileşen ve transfer öğrenme kurulumu açıklanmaktadır.

\begin{figure}[H]
\centering
\includegraphics[width=\linewidth]{architecture.pdf}
\caption{Çoklu-CNN özellik füzyon mimarisi: üç önceden eğitilmiş omurga (ResNet50,
MobileNetV2, EfficientNetB0; frozen ya da son 3 blok fine-tune, BatchNorm
dondurulmuş) $\to$ 512-boyutlu dal projeksiyonu (Linear$\to$LayerNorm$\to$GELU)
$\to$ füzyon (concat / weighted / GMU) $\to$ MLP. Soldaki kesikli kutu omurgaları,
ortadaki kesikli kutu projeksiyon dallarını gösterir. Çıkarımda final model TTA
(base/hflip/scale/scale\_hflip; softmax ortalaması) kullanır.}
\label{fig:arch}
\end{figure}

\subsection{Omurga Ağlar (Özellik Çıkarıcılar)}
\label{subsec:backbones}
Üç omurga ağ olarak ResNet50~\cite{he2016resnet},
MobileNetV2~\cite{sandler2018mobilenetv2} ve
EfficientNetB0~\cite{tan2019efficientnet} kullanılır. Her ağ ImageNet üzerinde
önceden eğitilmiş torchvision uygulamasından alınır ve sınıflandırma başlığı
kaldırılarak son havuzlama (pooling) katmanından özellik vektörü elde edilir.
Çıkış özellik boyutları ResNet50 için 2048, MobileNetV2 ve EfficientNetB0 için
1280'dir.
Tüm görüntüler ağlara $224\times224$ çözünürlükte ve ImageNet ortalama/standart
sapma değerleriyle normalize edilerek verilir; böylece üç omurga ağ kontrollü
ve karşılaştırılabilir koşullarda çalışır.

\subsection{Projeksiyon Katmanı}
\label{subsec:projection}
Omurga ağların özellik boyutları birbirinden farklı olduğundan (2048'e karşı
1280), füzyondan önce her dal ortak bir boyuta izdüşürülür. Her omurga ağın
özelliği 512 boyutlu bir projeksiyona eşlenir; bu projeksiyon kavramsal olarak
bir doğrusal katman, ardından katman normalizasyonu (LayerNorm) ve GELU
aktivasyonundan oluşur. Bu adım hem boyut uyumsuzluğunu giderir
hem de farklı ölçeklerdeki özellikleri füzyon için ortak bir temsile getirir.

\subsection{Birleştirme (Concatenation) Füzyonu}
\label{subsec:concat}
Birleştirme füzyonu, projeksiyon sonrası dal özelliklerini özellik ekseninde
uç uca ekleyerek tek bir vektör oluşturur. $N$ dal için her biri 512 boyutlu
olduğundan birleşik temsil $N\times512$ boyutludur (ör. üçlü birleşimde 1536).
Bu yöntem öğrenilebilir bir füzyon parametresi içermez ve tüm dalların bilgisini
eşit biçimde aşağı akışa taşır.

\subsection{Ağırlıklı (Weighted) Füzyon}
\label{subsec:weighted}
Ağırlıklı füzyon, dalların katkısını öğrenilebilir kılar. $N$ dal için
$\mathbf{a}\in\mathbb{R}^{N}$ öğrenilebilir bir parametre vektörüdür (uygulamada
sıfırlarla, yani düzgün başlangıçla ilklendirilir). Dal ağırlıkları bu vektörün
softmax'ı ile elde edilir ve projekte edilmiş dal özelliklerinin
$\mathbf{x}_i\in\mathbb{R}^{512}$ ağırlıklı toplamı alınır:
\begin{equation}
  \mathbf{w} = \operatorname{softmax}(\mathbf{a}), \qquad
  \mathbf{h} = \sum_{i=1}^{N} w_i\,\mathbf{x}_i ,
\end{equation}
burada $\sum_i w_i = 1$ ve birleşik temsil $\mathbf{h}$ yine 512 boyutludur.
Birleştirme füzyonundan farkı, hangi omurga ağın temsiline ne ölçüde
güvenileceğinin eğitim sırasında ayarlanabilmesidir.

\subsection{Gelişmiş Füzyon: GMU (İsteğe Bağlı Eksen)}
\label{subsec:gmu}
Zorunlu iki füzyon yöntemine ek olarak, gelişmiş bir kapılı çok-kipli birim
(Gated Multimodal Unit, GMU) varyantı da uygulanmış ve aynı protokol altında
değerlendirilmiştir~\cite{arevalo2017gmu}. Sabit ağırlıklı füzyondan farklı
olarak GMU, kapı (gate) ağırlıklarını her örnek için ayrı hesaplar. Bu projede
Arevalo ve ark.\ §3.1'deki iki-kipli formülasyonun $N$-dal softmax
genelleştirmesi kullanılır:
\begin{align}
  \mathbf{h}_i &= \tanh(W_i\,\mathbf{x}_i), \qquad i = 1,\dots,N \\
  \mathbf{z}   &= \operatorname{softmax}\!\big(W_z\,[\mathbf{x}_1,\dots,\mathbf{x}_N]\big) \in \mathbb{R}^{N} \\
  \mathbf{h}   &= \sum_{i=1}^{N} z_i\,\mathbf{h}_i ,
\end{align}
burada $[\cdot,\dots,\cdot]$ birleştirme işlemidir ve $W_z$ kapısı, dönüşümden
önceki ham dal özelliklerini girdi olarak alır (paper §3.1). Softmax, iki-kipli
$\sigma/(1-\sigma)$ çiftini, ağırlıkları toplamı 1 olan $N$-yönlü bir yarışmaya
genelleştirir.

\subsection{MLP Sınıflandırıcı}
\label{subsec:mlp}
Birleşik özellik vektörü, tek bir çok katmanlı algılayıcı (MLP) ile 23 sınıfa
sınıflandırılır. Sınıflandırıcı bir gizli katman (boyut 256) ve 0,3 oranında
dropout içerir. Sınıflandırıcı tüm deneylerde MLP olarak sabit tutulmuştur;
SVM, RandomForest veya XGBoost gibi alternatif sınıflandırıcılar
değerlendirilmemiştir. Bu sayede karşılaştırmalar omurga ağ,
füzyon ve transfer öğrenme eksenlerinde yoğunlaşır.

\subsection{Frozen Çıkarım ve Fine-tuning}
\label{subsec:transfer}
İki transfer öğrenme yaklaşımı karşılaştırılır:
\begin{itemize}
  \item \textbf{Frozen özellik çıkarımı:} Omurga ağların ağırlıkları
        dondurulur; özellikler bir kez çıkarılıp önbelleğe alınır ve yalnızca
        projeksiyon, füzyon ve MLP bileşenleri eğitilir.
  \item \textbf{Fine-tuning:} Her omurga ağın son üç bloğu çözülerek uçtan uca
        eğitilir. Bu kurulumda katman bazlı azalan öğrenme oranı (LLRD)
        kullanılır ve dondurulmuş bloklardaki BatchNorm istatistikleri,
        proje kuralı gereği güncellenmeden tutulur.
\end{itemize}
Her iki yaklaşımda da sınıf dengesizliği eğitim aşamasında ağırlıklı örnekleme
(weighted sampler) ile ele alınır. Eğitim ayrıntıları ve hiperparametreler
Bölüm~\ref{sec:setup}'te verilmiştir.

\subsection{Sınıf Dengesizliğine Duyarlı Kayıp: Focal Loss (Ablasyon)}
\label{subsec:focal}
Sınıf dengesizliğini kayıp düzeyinde ele almak için en iyi yapılandırma (exp 11,
üçlü ağırlıklı fine-tune) üzerinde tek değişkenli bir ablasyon olarak focal
loss~\cite{lin2017focal} denenmiştir. Lin ve ark.\ §3.2, eşitlik
(5)'teki $\alpha$-dengeli biçim kullanılır:
\begin{equation}
  \mathrm{FL}(p_t) = -\,\alpha_t\,(1 - p_t)^{\gamma}\,\log(p_t),
\end{equation}
burada $p_t$ gerçek sınıfın model olasılığı, $\gamma \ge 0$ odaklama
parametresidir; paper'ın önerdiği $\gamma = 2{,}0$ alınmıştır. $(1-p_t)^{\gamma}$
modülasyon çarpanı kolay örnekleri ($p_t > 0{,}5$) baskılayıp eğitimi zor
örneklere yöneltir. Bu ablasyon, focal loss ile etiket yumuşatmayı kasıtlı
olarak \emph{birleştirmez}: etiket yumuşatma hedef dağılımını değiştirir ve
$p_t$'ye dayanan focal modülasyonuyla çakışır. Bu nedenle focal yapılandırması,
exp 11'in CE yapılandırmasından yalnızca kayıp bloğunda ayrılır; diğer tüm
hiperparametreler aynıdır (temiz, tek değişkenli karşılaştırma).

\subsection{Test Zamanı Veri Artırma (TTA)}
\label{subsec:tta}
Nihai modelin çıkarım reçetesi, eğitim maliyeti olmadan uygulanan
deterministik bir test zamanı veri artırmadır (TTA). Her görüntü için dört
deterministik görünüm üretilir --- (i) \emph{base}: Resize(256)$\rightarrow$
CenterCrop(224); (ii) \emph{hflip}: base + yatay çevirme; (iii) \emph{scale}:
Resize(224)$\rightarrow$CenterCrop(224); (iv) \emph{scale\_hflip}: scale + yatay
çevirme --- ve bu görünümlerin softmax olasılıkları görüntü başına ortalanır. TTA
\emph{model başına} ve \emph{örnek başına} uygulanır; herhangi bir kat (fold)
sınırını aşmadığından sızıntıya yol açmaz.

\subsection{Sızıntısız Seed Topluluğu (Ensemble)}
\label{subsec:ensemble}
Resmi 5-katlı protokolde yalnızca sızıntısız topluluk yöntemine izin verilir.
Seed topluluğu, aynı yapılandırmanın farklı tohumlarla (seed)
eğitilmiş modellerinin softmax olasılıklarını \emph{her katın kendi ayrık test
kümesi içinde} ortalayıp ardından 5 katı havuzlayarak her örneğin tam olarak bir
kez görünmesini sağlar ($n = 10{.}662$). Olasılıklar her zaman softmax düzeyinde
ortalanır; argmax oylaması kullanılmaz. Katlar arası saf model ortalaması bu
çalışmada kullanılmamıştır: 5 kat modelini tüm veri kümesi üzerinde ortalamak
sızıntıya yol açar,
çünkü her örnek katların 4'ünün eğitim kümesinde yer almıştır. Bu nedenle
modeller hiçbir zaman katlar arasında topluluk haline getirilmez.
